% GENERATED FILE. Edit canonical lessons, then run npm run generate:books.
% canonical-source-hash: fnv1a64:ecd3ac2b0d9e4971
% canonical-lessons: HI-C358-seva, HI-C358-tyag, HI-C358-svarth, HI-C358-lalac, HI-C358-jatana, HI-R358-review-351-358

\chapter{Serving Others}
\label{ch:hi-serving-others}
\hlchaptermodality{358}

\begin{chapteropening}
\textbf{By the end of this chapter:} \emph{I can talk about selfless and selfish acts.}

Complete the last lesson of chapter 358: I can talk about selfless and selfish acts.
\end{chapteropening}

\section[\texorpdfstring{\hi{सेवा} (sevā)}{सेवा (sevā)}]{\hi{सेवा} (sevā) — service (to others)}

\label{lesson:HI-C358-seva}



\begin{quote}
Before the new one: say the Hindi for belief, then the Hindi for to adopt.
\end{quote}

\subsection*{You'll want to know: \hi{सेवा}}
\textbf{\hi{सेवा}} — \emph{sevā} — \textquotedblleft{}service (to others)\textquotedblright{}.

\begin{quote}
\hi{वह} \hi{बीमारों} \hi{की} \hi{सेवा} \hi{करती} \hi{है।}
\end{quote}

\begin{quote}
\emph{vah bīmāroṁ kī sevā kartī hai.}
\end{quote}

\begin{quote}
She looks after the sick.
\end{quote}

Feminine.

\begin{cousinweb}
\textbf{\hi{सेवा}} (\emph{sevā}) is a tatsama, borrowed from Sanskrit \emph{sevā}, \textquotedblleft{}service, attendance\textquotedblright{} (Monier-Williams).
\end{cousinweb}

\begin{grammarlens}[title={What you've built}]
One more word for giving reasons and opinions.
\end{grammarlens}

\subsection*{Your turn}
\begin{itemize}
\raggedright
  \item \emph{Say it:} \emph{sevā}
  \item \emph{Say it:} \emph{sevā}, once more
  \item \emph{Say it:} \emph{vah bīmāroṁ kī sevā kartī hai.} — She looks after the sick.
  \item \emph{Recall:} say the Hindi for belief, then the Hindi for to adopt, then say \emph{sevā} again
\end{itemize}

\begin{tcolorbox}[breakable,colback=teal!4,colframe=teal!35!black,title={Before you move on}]
What does \hi{सेवा} mean? (\textquotedblleft{}Service (to others)\textquotedblright{}.) Say it once more.
\end{tcolorbox}
\section[\texorpdfstring{\hi{त्याग} (tyāg)}{त्याग (tyāg)}]{\hi{त्याग} (tyāg) — renunciation, giving something up}

\label{lesson:HI-C358-tyag}



\begin{quote}
Before the new one: say the Hindi for to adopt, then the Hindi for service.
\end{quote}

\subsection*{You'll want to know: \hi{त्याग}}
\textbf{\hi{त्याग}} — \emph{tyāg} — \textquotedblleft{}renunciation, giving something up\textquotedblright{}.

\begin{quote}
\hi{माँ} \hi{ने} \hi{बच्चों} \hi{के} \hi{लिए} \hi{बहुत} \hi{त्याग} \hi{किया।}
\end{quote}

\begin{quote}
\emph{māṁ ne baccoṁ ke liye bahut tyāg kiyā.}
\end{quote}

\begin{quote}
Mother gave up a great deal for her children.
\end{quote}

Masculine.

\begin{cousinweb}
\textbf{\hi{त्याग}} (\emph{tyāg}) is a tatsama, borrowed from Sanskrit \emph{tyāga}, \textquotedblleft{}leaving, giving up\textquotedblright{} (Monier-Williams).
\end{cousinweb}

\begin{grammarlens}[title={What you've built}]
One more word for giving reasons and opinions.
\end{grammarlens}

\subsection*{Your turn}
\begin{itemize}
\raggedright
  \item \emph{Say it:} \emph{tyāg}
  \item \emph{Say it:} \emph{tyāg}, once more
  \item \emph{Say it:} \emph{māṁ ne baccoṁ ke liye bahut tyāg kiyā.} — Mother gave up a great deal for her children.
  \item \emph{Recall:} say the Hindi for to adopt, then the Hindi for service, then say \emph{tyāg} again
\end{itemize}

\begin{tcolorbox}[breakable,colback=teal!4,colframe=teal!35!black,title={Before you move on}]
What does \hi{त्याग} mean? (\textquotedblleft{}Renunciation, giving something up\textquotedblright{}.) Say it once more.
\end{tcolorbox}
\section[\texorpdfstring{\hi{स्वार्थ} (svārth)}{स्वार्थ (svārth)}]{\hi{स्वार्थ} (svārth) — self-interest, selfishness}

\label{lesson:HI-C358-svarth}



\begin{quote}
Before the new one: say the Hindi for service, then the Hindi for renunciation.
\end{quote}

\subsection*{You'll want to know: \hi{स्वार्थ}}
\textbf{\hi{स्वार्थ}} — \emph{svārth} — \textquotedblleft{}self-interest, selfishness\textquotedblright{}.

\begin{quote}
\hi{उसकी} \hi{हर} \hi{बात} \hi{में} \hi{स्वार्थ} \hi{होता} \hi{है।}
\end{quote}

\begin{quote}
\emph{uskī har bāt meṁ svārth hotā hai.}
\end{quote}

\begin{quote}
There is self-interest in everything he says.
\end{quote}

Masculine.

\begin{cousinweb}
\textbf{\hi{स्वार्थ}} (\emph{svārth}) is a tatsama, borrowed from Sanskrit \emph{svārtha}, \textquotedblleft{}one's own aim\textquotedblright{} (Monier-Williams).
\end{cousinweb}

\begin{grammarlens}[title={What you've built}]
One more word for giving reasons and opinions.
\end{grammarlens}

\subsection*{Your turn}
\begin{itemize}
\raggedright
  \item \emph{Say it:} \emph{svārth}
  \item \emph{Say it:} \emph{svārth}, once more
  \item \emph{Say it:} \emph{uskī har bāt meṁ svārth hotā hai.} — There is self-interest in everything he says.
  \item \emph{Recall:} say the Hindi for service, then the Hindi for renunciation, then say \emph{svārth} again
\end{itemize}

\begin{tcolorbox}[breakable,colback=teal!4,colframe=teal!35!black,title={Before you move on}]
What does \hi{स्वार्थ} mean? (\textquotedblleft{}Self-interest, selfishness\textquotedblright{}.) Say it once more.
\end{tcolorbox}
\section[\texorpdfstring{\hi{लालच} (lālac)}{लालच (lālac)}]{\hi{लालच} (lālac) — greed}

\label{lesson:HI-C358-lalac}



\begin{quote}
Before the new one: say the Hindi for renunciation, then the Hindi for self-interest.
\end{quote}

\subsection*{You'll want to know: \hi{लालच}}
\textbf{\hi{लालच}} — \emph{lālac} — \textquotedblleft{}greed\textquotedblright{}.

\begin{quote}
\hi{लालच} \hi{बुरी} \hi{चीज़} \hi{है।}
\end{quote}

\begin{quote}
\emph{lālac burī cīz hai.}
\end{quote}

\begin{quote}
Greed is a bad thing.
\end{quote}

Masculine.

\begin{cousinweb}
\textbf{\hi{लालच}} (\emph{lālac}) is a tadbhava, inherited from Sanskrit \emph{lālitya}, \textquotedblleft{}amorousness\textquotedblright{} (CDIAL 11029). CDIAL lists the Hindi \emph{lālac}, \textquotedblleft{}cupidity\textquotedblright{}, under this entry; how a word for amorousness came to mean greed is not spelled out there.
\end{cousinweb}

\begin{grammarlens}[title={What you've built}]
One more word for giving reasons and opinions.
\end{grammarlens}

\subsection*{Your turn}
\begin{itemize}
\raggedright
  \item \emph{Say it:} \emph{lālac}
  \item \emph{Say it:} \emph{lālac}, once more
  \item \emph{Say it:} \emph{lālac burī cīz hai.} — Greed is a bad thing.
  \item \emph{Recall:} say the Hindi for renunciation, then the Hindi for self-interest, then say \emph{lālac} again
\end{itemize}

\begin{tcolorbox}[breakable,colback=teal!4,colframe=teal!35!black,title={Before you move on}]
What does \hi{लालच} mean? (\textquotedblleft{}Greed\textquotedblright{}.) Say it once more.
\end{tcolorbox}
\section[\texorpdfstring{\hi{जताना} (jatānā)}{जताना (jatānā)}]{\hi{जताना} (jatānā) — to show, to make known (a feeling)}

\label{lesson:HI-C358-jatana}



\begin{quote}
Before the new one: say the Hindi for self-interest, then the Hindi for greed.
\end{quote}

\subsection*{You'll want to know: \hi{जताना}}
\textbf{\hi{जताना}} — \emph{jatānā} — \textquotedblleft{}to show, to make known (a feeling)\textquotedblright{}.

\begin{quote}
\hi{मीरा} \hi{ने} \hi{अपना} \hi{दुख} \hi{नहीं} \hi{जताया।}
\end{quote}

\begin{quote}
\emph{Mīrā ne apnā dukh nahīṃ jatāyā.}
\end{quote}

\begin{quote}
Mira didn't show her sorrow.
\end{quote}

Past, with \hi{ने} (he or she …): \hi{उसने} \hi{अपना} \hi{प्यार} \hi{जताया}, \emph{usne apnā pyār jatāyā}.

\begin{cousinweb}
\textbf{\hi{जताना}} (\emph{jatānā}) is a tadbhava, inherited from Sanskrit \emph{jñapta}, \textquotedblleft{}instructed\textquotedblright{} (CDIAL 5273). CDIAL glosses the Hindi verb \textquotedblleft{}to make known\textquotedblright{}.
\end{cousinweb}

\begin{grammarlens}[title={What you've built}]
One more word for giving reasons and opinions.
\end{grammarlens}

\subsection*{Your turn}
\begin{itemize}
\raggedright
  \item \emph{Say it:} \emph{jatānā}
  \item \emph{Say it:} \emph{jatānā}, once more
  \item \emph{Say it:} \emph{Mīrā ne apnā dukh nahīṃ jatāyā.} — Mira didn't show her sorrow.
  \item \emph{Recall:} say the Hindi for self-interest, then the Hindi for greed, then say \emph{jatānā} again
\end{itemize}

\begin{tcolorbox}[breakable,colback=teal!4,colframe=teal!35!black,title={Before you move on}]
What does \hi{जताना} mean? (\textquotedblleft{}To show, to make known (a feeling)\textquotedblright{}.) Say it once more.
\end{tcolorbox}
\section[(dialogue)]{Looking back: from Rights and Duties to Serving Others}

\label{lesson:HI-R358-review-351-358}



\begin{quote}
Say the Hindi for a right and a duty.
\end{quote}

\subsection*{Your turn}
Say these aloud:

\begin{itemize}
\raggedright
  \item a right, a duty, to abide by, proper, justice
  \item imprisonment, a judge, a verdict, guilty, bail
  \item parliament, governance, a policy, a proposal, to hold to be
  \item a vote, a citizen, public opinion, democracy, to impose
  \item a department, an application, an official, a command, to suggest
  \item equal, discrimination, bias, fairness, to provoke
  \item a value, integrity, honour, belief, to adopt
  \item service, renunciation, self-interest, greed, to make known
\end{itemize}

\begin{tcolorbox}[breakable,colback=teal!4,colframe=teal!35!black,title={Before you move on}]
A right? (\textbf{\hi{अधिकार}}, \emph{adhikār}.) A duty? (\textbf{\hi{कर्तव्य}}, \emph{kartavya}.) To abide by? (\textbf{\hi{पालन} \hi{करना}}, \emph{pālan karnā}.) Proper? (\textbf{\hi{उचित}}, \emph{ucit}.) Justice? (\textbf{\hi{न्याय}}, \emph{nyāy}.) Imprisonment? (\textbf{\hi{क़ैद}}, \emph{qaid}.) A judge? (\textbf{\hi{जज}}, \emph{jaj}.) A verdict? (\textbf{\hi{फ़ैसला}}, \emph{faislā}.) Guilty? (\textbf{\hi{दोषी}}, \emph{doṣī}.) Bail? (\textbf{\hi{ज़मानत}}, \emph{zamānat}.) Parliament? (\textbf{\hi{संसद}}, \emph{saṁsad}.) Governance? (\textbf{\hi{शासन}}, \emph{śāsan}.) A policy? (\textbf{\hi{नीति}}, \emph{nīti}.) A proposal? (\textbf{\hi{प्रस्ताव}}, \emph{prastāv}.) To hold to be? (\textbf{\hi{ठहराना}}, \emph{ṭhahrānā}.) A vote? (\textbf{\hi{वोट}}, \emph{voṭ}.) A citizen? (\textbf{\hi{नागरिक}}, \emph{nāgrik}.) Public opinion? (\textbf{\hi{जनमत}}, \emph{janmat}.) Democracy? (\textbf{\hi{लोकतंत्र}}, \emph{loktantra}.) To impose? (\textbf{\hi{थोपना}}, \emph{thopnā}.) A department? (\textbf{\hi{विभाग}}, \emph{vibhāg}.) An application? (\textbf{\hi{अर्ज़ी}}, \emph{arzī}.) An official? (\textbf{\hi{अधिकारी}}, \emph{adhikārī}.) A command? (\textbf{\hi{आदेश}}, \emph{ādeś}.) To suggest? (\textbf{\hi{सुझाना}}, \emph{sujhānā}.) Equal? (\textbf{\hi{बराबर}}, \emph{barābar}.) Discrimination? (\textbf{\hi{भेदभाव}}, \emph{bhedbhāv}.) Bias? (\textbf{\hi{पक्षपात}}, \emph{pakṣapāt}.) Fairness? (\textbf{\hi{इंसाफ़}}, \emph{iṁsāf}.) To provoke? (\textbf{\hi{उकसाना}}, \emph{uksānā}.) A value? (\textbf{\hi{मूल्य}}, \emph{mūlya}.) Integrity? (\textbf{\hi{ईमान}}, \emph{īmān}.) Honour? (\textbf{\hi{इज़्ज़त}}, \emph{izzat}.) Belief? (\textbf{\hi{आस्था}}, \emph{āsthā}.) To adopt? (\textbf{\hi{अपनाना}}, \emph{apnānā}.) Service? (\textbf{\hi{सेवा}}, \emph{sevā}.) Renunciation? (\textbf{\hi{त्याग}}, \emph{tyāg}.) Self-interest? (\textbf{\hi{स्वार्थ}}, \emph{svārth}.) Greed? (\textbf{\hi{लालच}}, \emph{lālac}.) To make known? (\textbf{\hi{जताना}}, \emph{jatānā}.)
\end{tcolorbox}
